\changecolours{ScoutGreen}
\section[Exercícios]{Exercícios e Desafios Práticos Resolvidos}

% ------------------------------------------------------------------------------
\twocolframe[title={Desafio 1: Tradução Tripla (Álgebra $\leftrightarrow$ CRT $\leftrightarrow$ SQL)},leftcol=7.0cm,rightcol=7.0cm]{%
  \alert{Enunciado do Problema}\par
  \itemseps[topsepi=0.5mm,itemsepi=1mm,labelsep=2mm,leftmargini=3mm]
  \begin{itemize}
    \item Recupere o \texttt{nome} e a \texttt{cidade} dos clientes de \texttt{'Tauá'} com compras individuais de valor superior a R\$ 400,00.
    \item Esquemas:
      \begin{itemize}
        \item \texttt{CLIENTE(id\_cli, nome, cidade)}
        \item \texttt{PEDIDO(id\_ped, id\_cli, valor)}
      \end{itemize}
    \item Formule a consulta nas 3 linguagens.
  \end{itemize}
}{%
  \alert{Resoluções Formais Equivalentes}\par
  \itemseps[topsepi=0.5mm,itemsepi=1mm,labelsep=2mm,leftmargini=3mm]
  \begin{itemize}
    \item \textbf{Álgebra Relacional:}
      {\tiny
      \[ \pi_{\text{nome, cidade}}\left(\sigma_{\text{cidade} = '\text{Tauá}' \,\land\, \text{valor} > 400}(\text{CLIENTE} \bowtie \text{PEDIDO})\right) \]}
    \item \textbf{Cálculo de Tuplas (CRT):}
      {\tiny
      \begin{align*}
        \{ t^{(2)} \mid &\exists c \in \text{CLIENTE} \, (c[\text{cidade}] = '\text{Tauá}' \land t[\text{nome}] = c[\text{nome}] \land t[\text{cid}] = c[\text{cid}] \\
        &\land \exists p \in \text{PEDIDO} \, (p[\text{id\_cli}] = c[\text{id\_cli}] \land p[\text{valor}] > 400)) \}
      \end{align*}}
    \item \textbf{SQL Padrão:}
      {\tiny\ttfamily
      SELECT c.nome, c.cidade\\
      FROM CLIENTE c JOIN PEDIDO p ON c.id\_cli = p.id\_cli\\
      WHERE c.cidade = 'Tauá' AND p.valor > 400;}
  \end{itemize}
}

% ------------------------------------------------------------------------------
\twocolframe[title={Desafio 2: Quantificação Universal e De Morgan no CRT},leftcol=7.0cm,rightcol=7.0cm]{%
  \alert{Enunciado do Problema}\par
  \itemseps[topsepi=0.5mm,itemsepi=1mm,labelsep=2mm,leftmargini=3mm]
  \begin{itemize}
    \item Encontre os clientes que compraram \textbf{todos} os livros promocionais com preço menor que R\$ 30,00.
    \item Formule em CRT utilizando o quantificador universal ($\forall$) e reescreva-o usando o quantificador existencial ($\exists$) via Leis de De Morgan.
  \end{itemize}
}{%
  \alert{Resoluções Lógicas}\par
  \itemseps[topsepi=0.5mm,itemsepi=1mm,labelsep=2mm,leftmargini=3mm]
  \begin{itemize}
    \item \textbf{Forma Universal ($\forall$):}
      {\tiny
      \begin{align*}
        \{ c[\text{nome}] \mid c \in \text{CLIENTE} \land \forall l \in \text{LIVRO} \, ( \\
        l[\text{preco}] < 30 \to \exists p \in \text{PEDIDO} \, ( \\
        p[\text{id\_cli}] = c[\text{id\_cli}] \land p[\text{isbn}] = l[\text{isbn}])) \}
      \end{align*}}
    \item \textbf{Forma Existencial ($\exists$ via De Morgan):}
      {\tiny
      \begin{align*}
        \{ c[\text{nome}] \mid c \in \text{CLIENTE} \land \neg \exists l \in \text{LIVRO} \, ( \\
        l[\text{preco}] < 30 \land \neg \exists p \in \text{PEDIDO} \, ( \\
        p[\text{id\_cli}] = c[\text{id\_cli}] \land p[\text{isbn}] = l[\text{isbn}])) \}
      \end{align*}}
    \item \emph{"Não existe nenhum livro barato que o cliente não tenha comprado"}.
  \end{itemize}
}

% ------------------------------------------------------------------------------
\twocolframe[title={Desafio 3: Diagnóstico e Correção de Insegurança},leftcol=7.0cm,rightcol=7.0cm]{%
  \alert{Expressão sob Julgamento}\par
  \itemseps[topsepi=0.5mm,itemsepi=1mm,labelsep=2mm,leftmargini=3mm]
  \begin{itemize}
    \item Avalie a seguinte consulta proposta por um estudante:
      \[ \psi = \{ t \mid \neg(t \in \text{PEDIDO}) \land t[\text{valor}] > 100 \} \]
    \item \textbf{Diagnóstico Técnico:}
      \begin{itemize}
        \item É \textbf{INSEGURA}.
        \item O conjunto de tuplas do universo que satisfazem a negação é infinito.
        \item O esquema da tupla resultante $t$ não está limitado a nenhuma tabela do catálogo.
      \end{itemize}
  \end{itemize}
}{%
  \alert{Formulação Segura Correta}\par
  \itemseps[topsepi=0.5mm,itemsepi=1mm,labelsep=2mm,leftmargini=3mm]
  \begin{itemize}
    \item Se o objetivo for: \emph{"Clientes que não realizaram nenhum pedido de valor > 100"}:
      {\footnotesize
      \begin{align*}
        \{ c[\text{id\_cli, nome}] \mid &c \in \text{CLIENTE} \land \\
        &\neg \exists p \in \text{PEDIDO} \, ( \\
        &p[\text{id\_cli}] = c[\text{id\_cli}] \land \\
        &p[\text{valor}] > 100) \}
      \end{align*}}
    \item Todos os valores de atributos do resultado provêm comprovadamente de $DOM(\psi)$.
  \end{itemize}
}

% ------------------------------------------------------------------------------
\twocolframe[title={Desafio 4: Análise de Custo de Refatoração de Árvores},leftcol=7.0cm,rightcol=7.0cm]{%
  \alert{Métricas da Loja Virtual}\par
  \itemseps[topsepi=0.5mm,itemsepi=1mm,labelsep=2mm,leftmargini=3mm]
  \begin{itemize}
    \item Tabela $R$ (Clientes): $|R| = 100.000$ tuplas.
    \item Tabela $S$ (Pedidos): $|S| = 1.000.000$ tuplas.
    \item Condição de filtro: $S.\text{valor} > 5000$ (apenas $200$ compras especiais).
    \item Consulta original não otimizada:
      \[ \pi_{\text{nome}}\left(\sigma_{S.\text{valor} > 5000 \,\land\, R.\text{id} = S.\text{id}}(R \times S)\right) \]
  \end{itemize}
}{%
  \alert{Comparativo Quantitativo de Custo}\par
  \itemseps[topsepi=0.5mm,itemsepi=1mm,labelsep=2mm,leftmargini=3mm]
  \begin{itemize}
    \item \textbf{Custo Não Otimizado:}
      \begin{itemize}
        \item $R \times S = 100.000 \times 1.000.000 = \mathbf{10^{11} \text{ tuplas}}$!
      \end{itemize}
    \item \textbf{Custo Otimizado (Push-Down Selection):}
      \begin{enumerate}
        \item Seleção prévia: $S' \leftarrow \sigma_{S.\text{valor} > 5000}(S) \implies \mathbf{200 \text{ tuplas}}$.
        \item Junção: $R \bowtie S' \implies \mathbf{200 \text{ tuplas}}$.
      \end{enumerate}
    \item \textbf{Fator de Otimização:} Redução de $10^{11} \to 200$ tuplas temporárias ($\mathbf{500.000.000\times}$ mais leve).
  \end{itemize}
}
